\documentclass[a4paper,11pt]{article}

%% -------------------------------------------------------
%%  Pacchetti
%% -------------------------------------------------------
\usepackage[T1]{fontenc}
\usepackage[utf8]{inputenc}
\usepackage[italian]{babel}
\usepackage{amsmath,amssymb}
\usepackage{booktabs}
\usepackage{siunitx}
\usepackage{geometry}
\usepackage{graphicx}
\usepackage{tikz}
\usepackage{circuitikz}
\usepackage{hyperref}
\usepackage{xcolor}
\usepackage{caption}
\usepackage{float}

\geometry{margin=25mm}

\sisetup{
  output-decimal-marker = {,},
  group-separator       = {\,},
  per-mode              = symbol
}

\hypersetup{
  colorlinks = true,
  linkcolor  = blue!60!black,
  urlcolor   = blue!60!black
}

%% -------------------------------------------------------
%%  Intestazione
%% -------------------------------------------------------
\title{\textbf{Calcolo dell'Impedenza Termica:\\
Piastra in Alluminio con Raffreddamento a Liquido}}
\author{Davide}
\date{\today}

%% -------------------------------------------------------
\begin{document}
\maketitle
\tableofcontents
\bigskip

%% =======================================================
\section{Introduzione e modello fisico}
%% =======================================================

L'impedenza termica $Z_{\mathrm{th}}$ (o, in regime stazionario, la resistenza termica
$R_{\mathrm{th}}$) descrive la caduta di temperatura per unità di potenza dissipata lungo
un percorso termico. Per un sistema lineare e stazionario:

\begin{equation}
  \Delta T = R_{\mathrm{th}}\, P
  \qquad\left[\si{\kelvin\per\watt}\right]
\end{equation}

dove $\Delta T$ è la differenza di temperatura tra la sorgente di calore e il riferimento
(in questo caso il fluido refrigerante), e $P$ è la potenza dissipata in watt.

Il sistema in esame è composto da:
\begin{enumerate}
  \item una \textbf{piastra in alluminio} di dimensioni
        $\SI{40}{mm}\times\SI{160}{mm}\times\SI{5}{mm}$
        (lo spessore di \SI{5}{mm} è lungo l'asse verticale del flusso di calore);
  \item un \textbf{flusso di acqua con glicole etilenico al 30\%} in volume
        sulla faccia inferiore della piastra, con portata volumetrica
        $\dot{V}=\SI{24}{l/min}$.
\end{enumerate}

Il flusso di calore è monodimensionale e verticale, dalla sorgente (superiore)
verso il fluido (inferiore). La resistenza totale è la serie di due contributi:

\begin{equation}
  \boxed{R_{\mathrm{th,tot}} = R_{\mathrm{cond}} + R_{\mathrm{conv}}}
\end{equation}

Il circuito analogico elettro-termico è rappresentato in Figura~\ref{fig:circuit}.

\begin{figure}[H]
\centering
\begin{tikzpicture}[font=\small]
  % Nodi
  \node[circle, fill=red!20, draw, minimum size=8mm] (Ts) at (0,0) {$T_s$};
  \node[circle, fill=orange!20, draw, minimum size=8mm] (Tc) at (4,0) {$T_c$};
  \node[circle, fill=blue!20, draw, minimum size=8mm] (Tf) at (8,0) {$T_f$};
  % Resistenze
  \draw[thick] (Ts) -- node[above]{$R_{\mathrm{cond}}$} (Tc);
  \draw[thick] (Tc) -- node[above]{$R_{\mathrm{conv}}$} (Tf);
  % Freccia potenza
  \draw[->, thick, red!60!black] (-1.5,0) -- node[above]{$P$} (Ts);
  % Etichette nodi
  \node[below=3mm] at (Ts) {sorgente};
  \node[below=3mm] at (Tc) {case};
  \node[below=3mm] at (Tf) {fluido};
\end{tikzpicture}
\caption{Schema circuito termico equivalente. $T_s$ = temperatura sorgente,
         $T_c$ = temperatura sulla faccia inferiore della piastra,
         $T_f$ = temperatura del fluido.}
\label{fig:circuit}
\end{figure}

%% =======================================================
\section{Geometria e dati di ingresso}
%% =======================================================

\begin{table}[H]
\centering
\caption{Dati geometrici della piastra.}
\begin{tabular}{llll}
\toprule
Simbolo & Grandezza & Valore & Unità \\
\midrule
$L$  & Lunghezza piastra         & 160   & \si{mm} \\
$W$  & Larghezza piastra         & 40    & \si{mm} \\
$t$  & Spessore piastra (asse~z) & 5     & \si{mm} \\
$A$  & Area superficiale ($L\times W$) & 64 & \si{cm^2} \\
\bottomrule
\end{tabular}
\end{table}

L'area della superficie attiva:
\begin{equation}
  A = L \times W = \SI{0.160}{m} \times \SI{0.040}{m} = \SI{6.4e-3}{m^2}
\end{equation}

%% =======================================================
\section{Resistenza termica di conduzione}
%% =======================================================

Il trasferimento di calore per conduzione stazionaria attraverso una parete piana
è governato dalla legge di Fourier:

\begin{equation}
  q = -k \frac{\mathrm{d}T}{\mathrm{d}z}
\end{equation}

dove $q\,[\si{W/m^2}]$ è il flusso di calore e $k\,[\si{W/(m\cdot K)}]$ è la
conducibilità termica del materiale. Per una geometria planare con spessore $t$,
area $A$ e potenza $P = q\cdot A$, la resistenza termica di conduzione è:

\begin{equation}
  \boxed{R_{\mathrm{cond}} = \frac{t}{k\,A}}
  \qquad \left[\si{K/W}\right]
\end{equation}

Per l'alluminio generico (serie 6xxx) si assume $k_{\mathrm{Al}} = \SI{160}{W/(m\cdot K)}$.

\begin{equation}
  R_{\mathrm{cond}} = \frac{\SI{5e-3}{m}}{\SI{160}{W/(m\cdot K)} \times \SI{6.4e-3}{m^2}}
  = \SI{4.88e-3}{K/W} \approx \SI{4.88}{mK/W}
\end{equation}

\noindent\textbf{Nota:} questo valore è molto piccolo rispetto alla resistenza convettiva,
a conferma che l'alluminio è un ottimo conduttore termico per questa geometria.

%% =======================================================
\section{Resistenza termica di convezione}
%% =======================================================

\subsection{Proprietà del fluido}

Il fluido refrigerante è acqua con glicole etilenico al 30\% in volume.
Le proprietà fisiche a temperatura operativa tipica ($\approx\SI{50}{\celsius}$) sono:

\begin{table}[H]
\centering
\caption{Proprietà termofisiche del fluido (EG 30\% vol., $T\approx\SI{50}{\celsius}$).}
\begin{tabular}{llll}
\toprule
Simbolo & Proprietà & Valore & Unità \\
\midrule
$\rho$  & Densità                  & 1040  & \si{kg/m^3}       \\
$c_p$   & Calore specifico         & 3800  & \si{J/(kg\cdot K)} \\
$\mu$   & Viscosità dinamica       & 0.82  & \si{mPa\cdot s}   \\
$k_f$   & Conducibilità termica    & 0.45  & \si{W/(m\cdot K)} \\
$\Pr$   & Numero di Prandtl        & 6.92  & --                \\
\bottomrule
\end{tabular}
\end{table}

Il numero di Prandtl è calcolato come:
\begin{equation}
  \Pr = \frac{\mu\, c_p}{k_f}
      = \frac{\SI{8.2e-4}{Pa\cdot s} \times \SI{3800}{J/(kg\cdot K)}}{\SI{0.45}{W/(m\cdot K)}}
      \approx 6.92
\end{equation}

\subsection{Portata e velocità media}

La portata volumetrica è:
\begin{equation}
  \dot{V} = \SI{24}{l/min} = \frac{24}{60\times 1000}\,\si{m^3/s}
          = \SI{4.0e-4}{m^3/s}
\end{equation}

Il canale di raffreddamento ha sezione rettangolare coincidente con la base della piastra
($h_c = \SI{5}{mm}$, $w_c = \SI{40}{mm}$):
\begin{equation}
  A_c = h_c \times w_c = \SI{5e-3}{m} \times \SI{0.040}{m} = \SI{2.0e-4}{m^2}
\end{equation}

La velocità media del fluido nel canale:
\begin{equation}
  \bar{v} = \frac{\dot{V}}{A_c} = \frac{\SI{4.0e-4}{m^3/s}}{\SI{2.0e-4}{m^2}}
          = \SI{2.0}{m/s}
\end{equation}

\subsection{Diametro idraulico}

Per un canale rettangolare con altezza $h_c$ e larghezza $w_c$:
\begin{equation}
  D_h = \frac{4\,A_c}{P_b} = \frac{4\,h_c\,w_c}{2\,(h_c+w_c)}
      = \frac{4\times\SI{5e-3}{m}\times\SI{0.040}{m}}{2\,(\SI{5e-3}{m}+\SI{0.040}{m})}
      \approx \SI{8.89}{mm}
\end{equation}

dove $P_b$ è il perimetro bagnato della sezione.

\subsection{Numero di Reynolds}

\begin{equation}
  \Rey = \frac{\rho\,\bar{v}\,D_h}{\mu}
       = \frac{1040 \times 2.0 \times \SI{8.89e-3}{}}{8.2\times10^{-4}}
       \approx 22\,547
\end{equation}

Poiché $\Rey \gg 4000$, il regime è \textbf{turbolento pienamente sviluppato}.

\subsection{Correlazioni per il numero di Nusselt}

\subsubsection{Correlazione di Dittus-Boelter}

Per fluido riscaldato ($n=0.4$), valida per $\Rey > 10^4$, $0.7 < \Pr < 160$,
$L/D_h > 60$:
\begin{equation}
  \Nu_{\mathrm{DB}} = 0.023\,\Rey^{0.8}\,\Pr^{0.4}
                    = 0.023 \times 22547^{0.8} \times 6.92^{0.4}
                    \approx 151.5
\end{equation}

\subsubsection{Correlazione di Gnielinski (consigliata)}

Più accurata nel range $3000 < \Rey < 5\times10^6$, $0.5 < \Pr < 2000$:

\begin{equation}
  f = \bigl(0.790\ln\Rey - 1.64\bigr)^{-2}
    = \bigl(0.790\times\ln 22547 - 1.64\bigr)^{-2}
    \approx 0.02537
\end{equation}

\begin{equation}
  \Nu_{\mathrm{G}} = \frac{(f/8)\,(\Rey-1000)\,\Pr}
                         {1 + 12.7\sqrt{f/8}\,(\Pr^{2/3}-1)}
                  = \frac{(0.02537/8)\times 21547 \times 6.92}
                         {1 + 12.7\sqrt{0.02537/8}\,(6.92^{2/3}-1)}
                  \approx 164.1
\end{equation}

\subsection{Coefficiente convettivo e resistenza}

Con la correlazione di Gnielinski (più conservativa e accurata):
\begin{equation}
  h_{\mathrm{conv}} = \frac{\Nu_{\mathrm{G}}\,k_f}{D_h}
                    = \frac{164.1 \times \SI{0.45}{W/(m\cdot K)}}{\SI{8.89e-3}{m}}
                    \approx \SI{8308}{W/(m^2\cdot K)}
\end{equation}

La resistenza termica convettiva riferita all'intera superficie $A$:
\begin{equation}
  \boxed{R_{\mathrm{conv}} = \frac{1}{h_{\mathrm{conv}}\,A}
                            = \frac{1}{\SI{8308}{W/(m^2\cdot K)}\times\SI{6.4e-3}{m^2}}
                            \approx \SI{18.8}{mK/W}}
\end{equation}

%% =======================================================
\section{Resistenza termica totale e risultati}
%% =======================================================

I due contributi sono in serie:
\begin{equation}
  R_{\mathrm{th,tot}} = R_{\mathrm{cond}} + R_{\mathrm{conv}}
                      = \SI{4.88}{mK/W} + \SI{18.8}{mK/W}
                      \approx \SI{23.7}{mK/W}
\end{equation}

\begin{table}[H]
\centering
\caption{Riepilogo dei contributi alla resistenza termica totale.}
\begin{tabular}{llcc}
\toprule
Contributo & Espressione & Valore [\si{mK/W}] & \% totale \\
\midrule
Conduzione alluminio & $t/(k_{\mathrm{Al}}\,A)$         & 4.88  & 20.6\% \\
Convezione (Gnielinski) & $1/(h_{\mathrm{conv}}\,A)$     & 18.8  & 79.4\% \\
\midrule
\textbf{Totale} & $R_{\mathrm{cond}}+R_{\mathrm{conv}}$ & \textbf{23.7} & 100\% \\
\bottomrule
\end{tabular}
\end{table}

La temperatura della sorgente termica rispetto al fluido vale quindi:
\begin{equation}
  \Delta T = R_{\mathrm{th,tot}}\times P
           \approx \SI{23.7e-3}{K/W} \times P
\end{equation}

Ad esempio, per una potenza $P = \SI{1000}{W}$:
\begin{equation}
  \Delta T \approx \SI{23.7e-3}{K/W} \times \SI{1000}{W} = \SI{23.7}{K}
\end{equation}

%% =======================================================
\section{Validità del modello e ipotesi}
%% =======================================================

\begin{enumerate}
  \item \textbf{Regime stazionario.} L'analisi fornisce la resistenza termica $R_{\mathrm{th}}$,
        non l'impedenza termica dinamica $Z_{\mathrm{th}}(t)$. Per i transitori è necessario
        includere la capacità termica $C_{\mathrm{th}} = \rho\,c_p\,V_{\mathrm{Al}}$
        e modellare il sistema con un circuito RC equivalente (modello di Foster o Cauer).

  \item \textbf{Flusso di calore monodimensionale.} Valido se la distribuzione di potenza
        sulla superficie è uniforme. In presenza di sorgenti puntiformi (es.\ chip IGBT),
        la resistenza di spreading aggiunge un contributo non trascurabile.

  \item \textbf{Canale a parete liscia.} La correlazione di Gnielinski si applica a canali
        a parete liscia. Strutture \textit{pin-fin} o corrugazioni migliorano
        significativamente $h_{\mathrm{conv}}$.

  \item \textbf{Proprietà del fluido costanti.} Le proprietà termofisiche del glicole
        variano con la temperatura; per analisi più accurate si dovrebbe iterare sulla
        temperatura media del film.

  \item \textbf{Resistenza di contatto.} Non è inclusa la resistenza di contatto
        tra la piastra e la struttura del canale ($R_{\mathrm{contact}}$), che in
        pratica dipende dalla rugosità superficiale e dall'eventuale strato di
        pasta termica.
\end{enumerate}

%% =======================================================
\section{Impedenza termica dinamica (schema RC)}
%% =======================================================

In regime dinamico, la risposta termica è modellata da una rete RC equivalente.
Per una singola massa termica (piastra in alluminio):

\begin{equation}
  C_{\mathrm{th}} = \rho_{\mathrm{Al}}\,c_{p,\mathrm{Al}}\,V_{\mathrm{Al}}
\end{equation}

Con $\rho_{\mathrm{Al}} = \SI{2700}{kg/m^3}$, $c_{p,\mathrm{Al}} = \SI{900}{J/(kg\cdot K)}$,
$V_{\mathrm{Al}} = L\times W\times t = \SI{3.2e-5}{m^3}$:

\begin{equation}
  C_{\mathrm{th}} = 2700 \times 900 \times \SI{3.2e-5}{m^3}
                  \approx \SI{77.8}{J/K}
\end{equation}

La costante di tempo termica associata alla conduzione:
\begin{equation}
  \tau = R_{\mathrm{cond}}\times C_{\mathrm{th}}
       = \SI{4.88e-3}{K/W} \times \SI{77.8}{J/K}
       \approx \SI{0.38}{s}
\end{equation}

L'impedenza termica nel dominio della frequenza (modello di Foster a un polo):
\begin{equation}
  Z_{\mathrm{th}}(j\omega) = \frac{R_{\mathrm{th,tot}}}{1 + j\omega\tau}
\end{equation}

Per pulse di potenza con durata $t_p$, la temperatura di picco si stima con
il fattore di forma dalla curva $Z_{\mathrm{th}}(t)$:
\begin{equation}
  \Delta T_{\mathrm{peak}} = Z_{\mathrm{th}}(t_p)\times P_{\mathrm{peak}}
\end{equation}

%% =======================================================
\section{Conclusioni}
%% =======================================================

La resistenza termica stazionaria del sistema piastra-alluminio con raffreddamento
a glicole è:

\begin{equation}
  R_{\mathrm{th,tot}} \approx \SI{23.7}{mK/W}
\end{equation}

Il contributo dominante (\SI{\approx 80}{\percent}) è la resistenza convettiva.
La resistenza di conduzione attraverso i \SI{5}{mm} di alluminio è trascurabile
(\SI{\approx 5}{mK/W}), confermando che l'alluminio non è il collo di bottiglia
termico. Per migliorare le prestazioni occorre agire sul lato convettivo:
aumentare la portata, ridurre il diametro idraulico con strutture \textit{pin-fin},
oppure utilizzare acqua pura (migliore conducibilità rispetto al glicole).

\end{document}
